\section*{Introduction}

Single-cell foundation models (\scfm{}s) such as
scGPT~\cite{cui2024scgpt}, Geneformer~\cite{theodoris2023geneformer},
scFoundation~\cite{hao2024scfoundation} and
MaxToki~\cite{maxtoki2026} are trained on millions of single-cell
transcriptomes and are reused for cell-type annotation, gene-network
inference and in-silico perturbation. Whether their internal
representations reflect regulatory biology, or mainly re-organise the
co-expression already present in the training data, cannot be settled by
benchmark accuracy
alone~\cite{kedzierska2023assessing,liu2023evaluating,wenteler2024perteval}.
Mechanistic interpretability asks this question directly: it studies what
a trained model's internal components represent and how they combine,
using probes, sparse autoencoders and causal
interventions~\cite{bricken2023monosemanticity,conmy2023acdc,wang2023ioi,sharkey2025open}.

Such analyses are long chains of code, statistics and interpretation, and
each link can fail. A null model can fail to randomise; an intervention
can edit the wrong tensor; a summary can describe a number as something
it is not. Large-language-model (LLM) agents can now carry out chains of
this length with little human
input~\cite{lu2024aiscientist,boiko2023autonomous,mbran2024chemcrow,yang2024sweagent,chan2024mlebench},
which makes the question of how to check their work practical. Agent
benchmarks for data-driven science and for reproducing published
analyses~\cite{chen2024scienceagentbench,siegel2024corebench,starace2025paperbench}
show that agents often produce plausible results that do not hold up on
inspection. For interpretability of biological models, a wrong result can
be confidently stated, internally consistent and hard for a reader to
detect.

This paper evaluates one answer to that problem. The framework has two
parts. First, each interpretability method is written as a markdown
\emph{specification} that an agent follows when running the analysis.
Each specification is written to a ten-section template, with methods,
parameters, validation criteria and known pitfalls. Seven of the eight
deployed specifications lacked the template's code-references section.
Second, the agent's outputs are checked with an \emph{audit checklist}
of ten methodological objections that peer reviewers raised repeatedly
on the author's earlier interpretability studies. We
deployed the framework on MaxToki-217M, a recently released
trajectory-aware \scfm{}~\cite{maxtoki2026}. An agent ran eight
interpretability pipelines over several working sessions. On the last
day, one of these sessions audited the outputs of all eight with the
checklist.

We then asked how reliable the resulting outputs were, and what the two
parts of the framework contribute. We did this in four ways.
(1)~\emph{Independent verification}: every number in the run summaries was searched for in
the stored run outputs by code, key statistics were re-derived, and the
intervention code was unit-tested; every error found was recorded in an
error ledger with the evidence that confirms it and with who or what found
it first. (2)~A \emph{blind re-audit} (Study~C): fresh agents audited a
copy of the project as it stood before its audit, under three
instructions (the checklist as deployed, a generic version of the
checklist, or a generic critical review of equal budget), and their
findings were scored against the ledger. (3)~A \emph{controlled
execution study} (Study~A): agents performed four analysis tasks built
from the deployment's data, with or without the specification. Each
output was kept as it was, and was also given a generic review and,
separately, a checklist review. After each review, a fresh agent repaired
the output, with the same repair instruction in both review arms. Graders
compared the final conclusions with fixed answer keys. They did not know
the review arm, but they could tell whether the agent had the
specification. (4)~A \emph{held-out benchmark} (Study~B): reviewers with or
without the checklist examined analyses from other projects and other
model families, each in a version with one documented error and in a
corrected version, so that missed errors and false alarms could both be
measured. Finally, we re-ran the main MaxToki-217M analyses with corrected code
and inputs and report what holds for that model. Some deployed analyses,
such as the topology screen and the ageing probe, were not re-run and are
not reported (Results).

\subsubsection*{What this study adds, and what it inherits}
The interpretability methods were developed in the author's earlier
studies of scGPT and
Geneformer~\cite{anon2026attention,anon2026spectral,anon2026topology,anon2026sae,anon2026circuits,anon2026manifold,anon2026exhaustive,anon2026ageing};
MaxToki-217M, the datasets and the agent models are third-party
resources. What was developed for this study is the framework itself (the
specifications and the checklist), its deployment on MaxToki-217M, the
verification and error ledger, the deterministic checks, the three
controlled studies and the corrected MaxToki-217M analyses
(Table~\ref{tab:inherited}). The three studies evaluate the framework;
the MaxToki-217M results are a case study and are stated for that model
only.

% TABLE: inherited (tab:inherited) is inserted here.
